\begin{definition}
In a Unique Games problem we are given a constraint graph $G = (V,E)$, a set of colors $[k] = \{1,\dotsc,k\}$ which is referred to as the alphabet, a set of variables $\{x_u\}_{u\in V}$, one for each vertex $u$, and a set of permutations (also known as constraints) $\pi_{uv} : [k]\to [k]$, one for each edge $uv\in E$. 
We study assignments giving a color from $[k]$ to each variable $x_u$, and are interested in the number of \emph{satisfied edges} (or satisfied constraints) of the form\footnote{Formally, let $G$ be an oriented graph so each edge has a direction. Then the edge (or constraint) $u\to v$ is satisfied when $\pi_{uv}(x_u)=x_v$, or equivalently $\pi_{uv}^{-1}(x_v)=x_u$. Without orientations it is unspecified whether to use $\pi_{uv}$ or $\pi_{uv}^{-1}$ here. We suppress this complication as in other parts of the paper it is more natural to consider undirected graphs.} $\pi_{uv}(x_u)=x_v$. The \emph{value} of the assignment is the fraction of satisfied constraints. 
The \emph{value} of the Unique Games instance is the maximum fraction of constraints that can be simultaneously satisfied.
%
We denote by \UG[k,\eps,\delta] the promise problem consisting of a Unique Games instance with alphabet size $k$ and the promise that the instance either has value at least $1-\eps$, or has value at most $\delta$. 
To solve the problem is to correctly determine which of the two cases hold. 
\end{definition}
%
When the parameters are unimportant or clear from context they are omitted, and we will always be interested in $\eps,\delta\ge 0$ with $1-\eps > \delta$, otherwise the problem is ill-posed or impossible to solve. 
The Unique Games Conjecture of Khot~\cite{Khot02a} can now be stated as follows.
%
\begin{conjecture*}[UGC]
For any constants $\eps,\delta > 0$ with $1-\eps>\delta$, there is a positive $k(\eps,\delta)$ such that for any alphabet size $k > k(\eps, \delta)$, the problem \UG[k,\eps,\delta] is NP-hard.
\end{conjecture*}
\noindent
We note that in~\cite{DKKMS18a,DKKMS18} it was shown that \UG[1/2-\delta,\delta] is NP-hard for any $\delta>0$.
%
Upon translating a Unique Games problem into a form amenable to methods from statistical physics, which we elaborate upon later, a natural variant of the problem arises.
%
\begin{definition}
Count Unique Games, or \CUG[f,k,\eps,\delta], is the promise problem consisting of a Unique Games instance with alphabet size $k$ and the promise that the instance either has at least $(fk)^{|V|}$ colorings with value at least $1-\eps$, or that every coloring has value at most $\delta$. 
To solve the problem is to correctly determine which of the two cases hold, and when the parameters are clear from context they are omitted. 
\end{definition}
%
Note that with $f=k^{-1}$, \CUG[f,k,\eps,\delta] corresponds exactly to \UG[k,\eps,\delta]; and with $f=1$ the problem is easily solvable as the guarantee covers all colorings in both cases. To solve \CUG[1,k,\eps,\delta] simply requires checking the value of an arbitrary coloring.
So via the parameter $f$ CUG can be smoothly reduced in difficulty from equal to UG to trivial. We can now state our main results.
%
